\documentclass[letterpaper,11pt]{moderncv}
\moderncvstyle{banking}
\moderncvcolor{black}

\usepackage[letterpaper, margin=0.8in]{geometry}
\usepackage{xcolor}
\usepackage{fontspec}
\setmainfont{Times New Roman}
\usepackage[unicode]{hyperref}

\name{Saeyoung}{Kim}
\address{Cambridge, MA, USA}{}{}
\phone[mobile]{+1 (351) 322 5510}
\email{saeyoung\_kim@uml.edu}
\social[linkedin]{saeyoung-macx-kim}
\social[github]{PrimaryAnomaly}

\recipient{Hiring Committee}{Space Sciences Laboratory\\University of California, Berkeley\\Berkeley, CA}
\date{\today}
\opening{Dear Hiring Committee,}
\closing{Sincerely,}

\begin{document}

\makelettertitle

I am writing to apply for the Safety and Mission Assurance Engineer position at the Space Sciences Laboratory. My background combines satellite integration and validation, embedded instrument development, and hands-on systems engineering across university and industry environments. The opportunity to support Berkeley SSL's spaceflight programs is especially compelling to me because it brings together spacecraft hardware, instrument-level engineering, and the disciplined testing and coordination required for successful missions.

At Hanwha Systems, I worked as a Systems Integration Engineer on military SAR satellite programs, where I supported electrical integration and validation across complex test and assembly campaigns. I developed and executed Electric Test Bed procedures, defined test plans for embedded and communication subsystems, contributed to Electrical Ground Support Equipment architecture, and diagnosed hardware-software issues during integration activities. This experience strengthened my ability to work methodically in aerospace environments where interoperability, traceability, and rigorous validation are essential.

My doctoral research in Electrical Engineering further developed my space systems foundation. I designed and fabricated multi-modal electrochemical and optical sensor hardware for CubeSat payload applications, developed MCU firmware in C/C++, and built embedded prototypes from circuit design and PCB layout through assembly and validation. That work required close coordination across disciplines, careful technical documentation, and sustained attention to system reliability at the instrument level.

In my current postdoctoral role at the University of Massachusetts Lowell, I develop technical software and model-based workflows for research use, which has further strengthened my documentation, cross-functional collaboration, and engineering analysis skills. Across these roles, I have consistently worked at the interface of hardware, software, testing, and research execution.

I would welcome the opportunity to bring my experience in satellite integration, embedded systems, instrument development, and technical collaboration to Berkeley SSL. Thank you for your time and consideration.

\makeletterclosing

\end{document}
